 
\chapter{Disentangle Mixture Distributions and Instrumental Variable Inequalities}
 \label{chapter::iv-inequalities}

 
 \cref{chapter::iv-experiment} 中的 IV model 施加了 \cref{assume::random}--\cref{assume::ER}：
 \begin{enumerate}
 \item
 $Z\ind \{D(1), D(0) ,Y(1), Y(0)   \}$;
 \item
 $\pr(U = \df) = 0$;
 \item 
 当 $U=\at$ 或 $\nt$ 时，$Y(1) = Y(0)$。
 \end{enumerate}


\cref{table::obs-latent} 总结了 monotonicity assumption 下的 observed groups 以及相应的 latent groups。


\begin{table}[h]
\centering
\caption{\cref{assume::Mono} 下的 observed groups 与 latent groups}\label{table::obs-latent}
\begin{tabular}{cccc}
\hline 
$Z$ & $D$ & $D(1)$ & latent groups \\
\hline 
$Z=1$ & $D=1$ & $D(1) = 1$ & $U = \cp$ 或 $\at$ \\
$Z=1$ & $D=0$ & $D(1) = 0$ & $U = \nt$ \\ 
$Z=0$ & $D=1$ & $D(0) = 1$ & $U = \at$  \\ 
$Z=0$ & $D=0$ & $D(0) = 0$ & $U = \cp$ 或 $\nt$ \\
\hline 
\end{tabular}
\end{table}


有意思的是，\cref{assume::random}--\cref{assume::ER} 合在一起会给出一些可检验的含义。\citet{balke1997bounds} 将它们称为 {\it instrumental variable inequalities}。本章给出这些 inequalities 一个特殊情形的直观推导。证明的核心很直接：先识别由 $U$ 定义的各个 latent groups 中 potential outcomes 的均值。

\section{Disentangle Mixture Distributions}
\label{sec::disentangle-mixtures}


我们先把主要结果总结为下面的 \cref{thm::cace-mixture-components}。定义
$$
\pi_u = \pr(U = u) \quad  \quad  (u=\at, \nt, \cp) 
$$
为类型 $U=u$ 的比例，并定义
$$
\mu_{zu} = E\{  Y(z) \mid U = u\}\quad  \quad (z=0,1;u=\at, \nt, \cp) .
$$ 
为类型 $U=u$ 的 potential outcome $Y(z)$ 的均值。Exclusion restriction 意味着 $\mu_{1\nt} = \mu_{0\nt} $ 且 $\mu_{1\at} = \mu_{0\at} $。分别用 $\mu_{\nt}$ 和 $\mu_{\at} $ 表示这两个共同值。

\begin{theorem}
\label{thm::cace-mixture-components}
在 \cref{assume::random}--\cref{assume::ER} 下，latent types 的比例可以由下式识别：
\begin{eqnarray*}
 \pi_\nt &=& \pr(D=0\mid Z=1)  ,\\
  \pi_\at &=&\pr(D=1\mid Z=0) , \\
  \pi_\cp &=& E(D\mid Z=1) - E(D\mid Z=0) ,
\end{eqnarray*}
而 potential outcomes 的 type-specific means 可以由下式识别：
\begin{eqnarray*}
  \mu_{\nt} &=& E(Y\mid Z=1, D=0) ,\\
 \mu_{\at} &=& E(Y\mid Z=0, D=1),\\
 \mu_{1\cp} &=& \pi_\cp^{-1}\left\{   E(DY\mid Z=1 )  -  E(DY\mid Z=0)    \right\}, \\
  \mu_{0\cp}  &=& \pi_\cp^{-1} \big[   E\{ (1-D)Y\mid Z=0 \}  -  E\{(1-D)Y\mid Z=1 \}    \big]. 
\end{eqnarray*}
\end{theorem}
 
 
\begin{myproof}{Theorem}{\cref{thm::cace-mixture-components}}
第一部分：我们先识别 latent compliance types 的比例。
Never takers 的比例可由下式识别：
\begin{eqnarray*}
\pr(D=0\mid Z=1) 
%&=& \pr\{ D(1) = 0\mid Z=1 \} = \pr\{ D(1) = 0, D(0) = 0 \mid Z=1 \} \\
&=& \pr(U=\nt \mid Z=1)  \\
&=& \pr(U=\nt) \\
&=&  \pi_\nt ,
\end{eqnarray*}
Always takers 的比例可由下式识别：
\begin{eqnarray*}
\pr(D=1\mid Z=0) 
%&=& \pr\{  D(0) = 1\mid Z=0 \} =  \pr\{ D(1)=1,  D(0) = 1\mid Z=0 \} \\
&=& \pr(U=\at \mid Z=0)  \\
&=& \pr(U=\at) =  \pi_\at . 
\end{eqnarray*}
因此，compliers 的比例为
\begin{eqnarray*}
\pi_\cp &=& \pr(U=\cp) \\
&=& 1-  \pi_\nt -  \pi_\at  \\
&=& 1 - \pr(D=0\mid Z=1)  - \pr(D=1\mid Z=0) \\
%&=&\pr(D=1\mid Z=1)  - \pr(D=1\mid Z=0) 
&=& E(D\mid Z=1) - E(D\mid Z=0) \\
&=& \tau_D,
\end{eqnarray*}
这与前面的讨论一致。虽然我们无法知道每个个体的 latent compliance type，但可以识别 never-takers、always-takers 和 compliers 的总体比例。


第二部分：接下来识别各 latent compliance type 内 potential outcomes 的均值。
%Under \cref{assume::ER}, 
%$$
%\mu_{1\at} = \mu_{0\at} \equiv \mu_{\at},\quad \mu_{1\nt} = \mu_{0\nt} \equiv \mu_{\nt} .
%$$
观测组 $(Z=1,D=0)$ 中只有 never takers，因此
$$
E(Y\mid Z=1, D=0) = E\{  Y(1) \mid Z=1, U=\nt \} = E\{  Y(1) \mid   U=\nt \}  =
\mu_{\nt}  . 
$$
观测组 $(Z=0,D=1)$ 中只有 always takers，因此
$$
E(Y\mid Z=0, D=1) = E\{  Y(0) \mid Z=0, U=\at \} = E\{  Y(0) \mid   U=\at \}  =
\mu_{\at} . 
$$
观测组 $(Z=1,D=1)$ 同时包含 compliers 和 always takers，因此
\begin{eqnarray*}
E(Y\mid Z=1, D=1) &=& E\{ Y(1)\mid Z=1, D(1)=1\} \\
&=& E\{ Y(1)\mid D(1)=1\}  \\   
&=& \pr\{  D(0)=1 \mid D(1)=1 \} E\{ Y(1)\mid D(1)=1, D(0)=1\}  \\
&&+ \pr\{  D(0)=0 \mid D(1)=1 \}E\{ Y(1)\mid D(1)=1, D(0)=0\}   \\
&=& \frac{\pi_\cp}{\pi_\cp + \pi_\at} \mu_{1\cp} +  \frac{\pi_\at}{\pi_\cp + \pi_\at} \mu_{\at}.
\end{eqnarray*}
解上面的线性方程可得
\begin{eqnarray*}
\mu_{1\cp} &=& \pi_\cp^{-1} \left\{  (\pi_\cp + \pi_\at)  E(Y\mid Z=1, D=1)  - \pi_\at    E(Y\mid Z=0, D=1)    \right\} \\
&=& \pi_\cp^{-1} \left\{  \pr(D=1\mid Z=1)  E(Y\mid Z=1, D=1) \right.\\
&&\left. \qquad  -\pr(D=1\mid Z=0) E(Y\mid Z=0, D=1)    \right\} \\
&=&  \pi_\cp^{-1}\left\{   E(DY\mid Z=1 )  -  E(DY\mid Z=0)    \right\}.
\end{eqnarray*}
观测组 $(Z=0,D=0)$ 同时包含 compliers 和 never takers，因此
\begin{eqnarray*}
E(Y\mid Z=0,D=0) &=& E\{ Y(0)\mid Z=0, D(0)=0\} \\
&=& E\{ Y(0)\mid D(0)=0\} \\ 
&=& \pr\{  D(1)=1 \mid D(0)=0 \} E\{ Y(0)\mid D(1)=1, D(0)=0\}  \\
&&+ \pr\{  D(1)=0 \mid D(0)= 0\}E\{ Y(0)\mid D(1)=0, D(0)=0\}   \\
&=& \frac{\pi_\cp}{\pi_\cp + \pi_\nt} \mu_{0\cp} +  \frac{\pi_\nt}{\pi_\cp + \pi_\nt} \mu_{\nt}.
\end{eqnarray*}
解上面的线性方程可得
\begin{eqnarray*}
\mu_{0\cp}  &=& \pi_\cp^{-1} \left\{  (\pi_\cp + \pi_\nt)  E(Y\mid Z=0, D=0)  - \pi_\nt    E(Y\mid Z=1, D=0)    \right\} \\
&=& \pi_\cp^{-1} \left\{ \pr(D=0\mid Z=0)  E(Y\mid Z=0, D=0)  \right.\\
&&\left. \qquad  -\pr(D=0\mid Z=1)   E(Y\mid Z=1, D=0)    \right\} \\
&=& \pi_\cp^{-1} \big[   E\{ (1-D)Y\mid Z=0 \}  -  E\{(1-D)Y\mid Z=1 \}    \big]. 
\end{eqnarray*}
\end{myproof}

利用 \cref{thm::cace-mixture-components} 中 $\mu_{1\cp} $ 和 $\mu_{0\cp} $ 的公式，可以将 $\tau_\cp = \mu_{1\cp}  - \mu_{0\cp}  $ 简化为
$$
\tau_\cp =  \{  E(Y\mid Z=1) - E(Y\mid Z=0) \}/\pi_\cp ,
$$
这正是前面 \cref{thm::CACE-identification} 中的公式。

\cref{thm::cace-mixture-components} 关注的是识别 potential outcomes 的均值 $\mu_{zu} $。\citet{imbens1997estimating} 推导了 potential outcomes 分布的更一般 identification formulas；细节留给 \cref{hw::disentangle-mixture-distrubution}。



\section{Testable implications: Instrumental Variable Inequalities}


通过 \cref{thm::cace-mixture-components} 绕一圈来推导 $\tau_\cp$ 的公式，是否还有额外价值？答案是肯定的。对于 binary outcome，下面的 inequalities 必须成立：
$$
0\leq \mu_{1\cp} \leq 1, \quad 
0\leq \mu_{0\cp} \leq 1, 
$$
这蕴含四个 inequalities：
\begin{eqnarray*}
E(DY\mid Z=1 )  -  E(DY\mid Z=0) &\geq & 0 ,\\
E(DY\mid Z=1 )  -  E(DY\mid Z=0) &\leq  & E(D\mid Z=1) - E(D\mid Z=0)  ,\\ 
E\{ (1-D)Y\mid Z=0 \}  -  E\{(1-D)Y\mid Z=1 \}   &\geq & 0,\\
E\{ (1-D)Y\mid Z=0 \}  -  E\{(1-D)Y\mid Z=1 \}  &\leq & E(D\mid Z=1) - E(D\mid Z=0) . 
\end{eqnarray*}
整理各项后，可以得到下面统一形式的 inequalities。


\begin{theorem}[Instrumental Variable Inequalities]\label{thm::iv-inequality}
若 outcome $Y$ 为 binary，则 \cref{assume::random}--\cref{assume::ER} 蕴含
\begin{eqnarray}\label{eq::iv-inequalities}
E(Q\mid Z=1) - E(Q\mid Z=0) \geq 0,  
\end{eqnarray}
其中 $Q = DY, D(1-Y), (D - 1)Y$ 以及 $D+Y-DY$。
\end{theorem}


在 IV assumptions \cref{assume::random}--\cref{assume::ER} 下，对于 $Q = DY, D(1-Y), (D - 1)Y$ 和 $D+Y-DY$，处理组与对照组的均值差都必须非负。重要的是，这些含义只涉及观测变量的分布。因此，拒绝 IV inequalities 就会拒绝 IV assumptions。


\citet{balke1997bounds} 在假设和不假设 monotonicity 的两种情形下推导了更一般的 IV inequalities。上面的证明策略来自 \citet{jiang2020measurement}，他们处理的是稍复杂一些的设定。\cref{thm::iv-inequality} 只陈述了 binary outcome 下的 testable implications。\cref{hw::ivineq-binary} 给出等价形式，\cref{hw::ivineq-general-Y} 给出 general outcome 下的结果。



\section{Examples}\label{section::binary-y-ivineq}


对于 binary outcome，可以用下面的 method of moments 估计所有参数。


\begin{lstlisting}
IVbinary = function(n111, n110, n101, n100, 
                    n011, n010, n001, n000){
  
  n_tr = n111 + n110 + n101 + n100
  n_co = n011 + n010 + n001 + n000
  n    = n_tr + n_co
  
  ## proportions of the latent strata
  pi_n = (n101 + n100)/n_tr
  pi_a = (n011 + n010)/n_co
  pi_c = 1 - pi_n - pi_a
  
  ## four observed means of the outcomes (Z=z,D=d)
  mean_y_11 = n111/(n111 + n110)
  mean_y_10 = n101/(n101 + n100)
  mean_y_01 = n011/(n011 + n010)
  mean_y_00 = n001/(n001 + n000)
  
  ## means of the outcomes of two strata
  mu_n1 = mean_y_10
  mu_a0 = mean_y_01
  ## ER implies the following two means
  mu_n0 = mu_n1
  mu_a1 = mu_a0
  ## stratum (Z=1,D=1) is a mixture of c and a
  mu_c1 = ((pi_c + pi_a)*mean_y_11 - pi_a*mu_a1)/pi_c
  ## stratum (Z=0,D=0) is a mixture of c and n
  mu_c0 = ((pi_c + pi_n)*mean_y_00 - pi_n*mu_n0)/pi_c
  
  ## identifiable quantities from the observed data
  list(pi_c = pi_c, 
       pi_n = pi_n, 
       pi_a = pi_a, 
       mu_c1= mu_c1,
       mu_c0= mu_c0,
       mu_n1= mu_n1,
       mu_n0= mu_n0,
       mu_a1= mu_a1,
       mu_a0= mu_a0,
       tau_c= mu_c1 - mu_c0)
}
\end{lstlisting}


下面重新考察两个 binary data 的经典例子。

\begin{example}\label{eg::1-iv-inequalities}
\citet{improve2014endovascular} 评估急诊腔内修复（emergency endovascular）与开放手术修复（open surgical repair）两种策略对临床诊断为破裂性主动脉瘤（ruptured aortic aneurism）患者的有效性。患者被随机分配到 endovascular strategy 或 open repair strategy。主要结局（primary outcome）是 30 天后的生存状态（survival status）。令 $Z$ 表示 treatment assigned，其中 $Z=1$ 对应 endovascular strategy，$Z=0$ 对应 open repair。令 $D$ 表示 treatment received。
令 $Y$ 表示 survival status，其中 $Y=1$ 表示死亡，$Y=0$ 表示存活。\cref{tab::data1} 总结了 observed data。
%The estimate of $\tau_\cp $ is $0.131$ with 95\% confidence interval $(-0.036, 0.298)$ including $0$.
使用上面的 \ri{IVbinary} 函数，可以得到如下估计：
\begin{lstlisting}
> investigators_analysis = IVbinary(n111 = 107,
+                                   n110 = 42,
+                                   n101 = 68,
+                                   n100 = 42,
+                                   n011 = 24,
+                                   n010 = 8,
+                                   n001 = 131,
+                                   n000 = 79)
>                         
> investigators_analysis
$pi_c
[1] 0.4430582

$pi_n
[1] 0.4247104

$pi_a
[1] 0.1322314

$mu_c1
[1] 0.7086064

$mu_c0
[1] 0.6292042

$mu_n1
[1] 0.6181818

$mu_n0
[1] 0.6181818

$mu_a1
[1] 0.75

$mu_a0
[1] 0.75

$tau_c
[1] 0.07940223
\end{lstlisting}
没有证据表明 IV assumptions 被违反。
\end{example}


\begin{example}\label{eg::2-iv-inequalities}
在 \citet{hirano2000assessing} 中，医生（physicians）被随机抽取，并收到一封鼓励其为 flu 高风险患者接种疫苗的信。实际接种 flu shot 是 treatment，outcome 是 flu-related hospital visits 的指示变量。不过，有些患者并不遵从其 assignment。令 $Z_i$ 表示是否受到接种 flu shot 的 encouragement：若 physician 收到 encouragement letter，则 $Z=1$，否则 $Z=0$。令 $D$ 表示 treatment received。
令 $Y$ 表示 outcome，其中若冬季发生 flu-related hospitalization 则 $Y=0$，否则 $Y=1$。数据的更多细节见 \cref{problem::flushot}。\cref{tab::data2} 总结了 observed data。
%The estimate of $\tau_\cp $ is $0.116$ with 95\% confidence interval $(-0.061, 0.293)$ including $0$.
使用上面的 \ri{IVbinary} 函数，可以得到如下估计：
\begin{lstlisting}
> flu_analysis = IVbinary(n111 = 31,
+                         n110 = 422,
+                         n101 = 84,
+                         n100 = 935,
+                         n011 = 30,
+                         n010 = 233,
+                         n001 = 99,
+                         n000 = 1027)
> flu_analysis
$pi_c
[1] 0.1183997

$pi_n
[1] 0.6922554

$pi_a
[1] 0.1893449

$mu_c1
[1] -0.004548064

$mu_c0
[1] 0.1200094

$mu_n1
[1] 0.08243376

$mu_n0
[1] 0.08243376

$mu_a1
[1] 0.1140684

$mu_a0
[1] 0.1140684

$tau_c
[1] -0.1245575
\end{lstlisting}
由于 $\hat\mu_{1\cp} < 0$，这里有证据表明 IV assumptions 被违反。
\end{example} 




 \begin{table}[t]
\caption{Binary data 与 IV inequalities}
\begin{center}
\begin{subtable}{0.5\textwidth}
\caption{\citet{improve2014endovascular} 的研究}
\begin{tabular}{cccccc} 
\hline 
      &   \multicolumn{2}{c}{$Z=1$}  &  &  \multicolumn{2}{c}{$Z=0$}  \\  \cline{2-3} \cline{5-6} 
      & $D=1$  &  $D=0$           &       &  $D=1$  &  $D=0$\\
 $Y=1$ &   107    &      68         &      &    24       &    131        \\
 $Y=0$ &   42    &     42            &      &     8         &     79           \\  
 \hline
\end{tabular}
\label{tab::data1}
\end{subtable}% 
\vspace{0.1in}
\begin{subtable}{0.5\textwidth}
\caption{\citet{hirano2000assessing} 的研究}
\begin{tabular}{cccccc}
\hline 
      &   \multicolumn{2}{c}{$Z=1$}  &  &  \multicolumn{2}{c}{$Z=0$}  \\  \cline{2-3} \cline{5-6} 
      & $D=1$  &  $D=0$           &       &  $D=1$  &  $D=0$\\
 $Y=1$ &   31    &     85         &      &     30       &    99       \\
 $Y=0$ &   424   &     944            &      &     237         &    1041           \\  
 \hline
\end{tabular}
\label{tab::data2}
\end{subtable}
\end{center}
\end{table}%

 
 
\section{Homework problems}

\homeworksolutionref{22}
\paragraph{More detailed data analysis}\label{hw::more-data-binary}

\cref{eg::1-iv-inequalities,eg::2-iv-inequalities} 忽略了估计中的 uncertainty。请计算真参数的 confidence intervals。
 



\paragraph{Risk ratio for compliers}\label{hw::rr-compliers}


对于 binary outcome，可以将 compliers 的 risk ratio 定义为
$$
\textsc{rr}_\cp = \frac{  \pr\{  Y(1) = 1\mid U = \cp \}  }{ \pr\{  Y(0) = 1\mid U = \cp \}   } .
$$

证明：在 \cref{assume::random}--\cref{assume::ER} 下，它可以由下式识别：
$$
\textsc{rr}_\cp =  \frac{  E(DY\mid Z=1)  - E(DY\mid Z=0)  }{  E\{  (D-1)Y\mid Z=1 \}   - E\{ (D-1)Y\mid Z=0 \} }.
$$



注：利用 \cref{thm::cace-mixture-components}，可以识别 $E\{  Y(1) \mid U = \cp \} $ 与 $E\{  Y(0) \mid U = \cp \} $ 之间的任意比较。


\paragraph{Disentangle the mixtures: distributional results}\label{hw::disentangle-mixture-distrubution}


本题推广 \cref{thm::cace-mixture-components}。定义
$$
f_{zu}(y) = \pr\{  Y(z) = y  \mid U=u \},  \quad  (z=0,1;u=\at, \nt, \cp)
$$
为 latent stratum $U=u$ 中 $Y(z)$ 的 density，并定义
$$
g_{zd}(y) = \pr(Y = y\mid Z=z, D=d)
$$
为 observed group $(Z=z, D=d)$ 内 outcome 的 density。Exclusion restriction 意味着 $f_{1\nt}(y) = f_{0\nt} (y) $ 且 $f_{1\at}(y) = f_{0\at}(y) $。分别用 $f_{\nt}(y)$ 和 $ f_{\at} (y) $ 表示这两个共同 density。



证明下面的 \cref{thm::cace-mixture-components-distribution}。


 
\begin{theorem}
\label{thm::cace-mixture-components-distribution}
在 \cref{assume::random}--\cref{assume::ER} 下，potential outcomes 的 type-specific densities 可以由下式识别：
\begin{eqnarray*}
  f_{\nt}(y) &=& g_{10}(y) ,\\
  f_{\at} (y) &=& g_{01}(y),\\
f_{1\cp}(y) &=& \pi_\cp^{-1} \{   \pr(D=1\mid Z=1 )g_{11}(y)  -  \pr(D=1\mid Z=0) g_{01}(y)   \} , \\
f_{0\cp} (y) &=& \pi_\cp^{-1} \{  \pr(D=0\mid Z=0)  g_{00}(y)  -  \pr(D=0\mid Z=1)  g_{10}(y)    \} . 
\end{eqnarray*}
\end{theorem}



 





\paragraph{Alternative form of \cref{thm::iv-inequality}}
\label{hw::ivineq-binary}

\cref{eq::iv-inequalities} 中的 inequalities 可以改写为
\begin{eqnarray*}
\pr(D=1, Y=y\mid Z=1) &\geq & \pr(D=1, Y=y\mid Z=0), \\
\pr(D=0, Y=y\mid Z=0) &\geq & \pr(D=0, Y=y\mid Z=1)
\end{eqnarray*}
对 $y=0,1$ 都成立。



\paragraph{IV inequalities for a general outcome}
\label{hw::ivineq-general-Y}

对于 general outcome $Y$，证明 \cref{assume::random}--\cref{assume::ER} 蕴含
\begin{eqnarray*}
\pr(D=1, Y\geq y\mid Z=1) &\geq & \pr(D=1, Y\geq y\mid Z=0), \\
\pr(D=1, Y <  y\mid Z=1) &\geq & \pr(D=1, Y <  y\mid Z=0), \\
\pr(D=0, Y\geq y\mid Z=0) &\geq & \pr(D=0, Y\geq y\mid Z=1),\\
\pr(D=0, Y <  y\mid Z=0) &\geq & \pr(D=0, Y <  y\mid Z=1)
\end{eqnarray*}
对所有 $y$ 都成立。
 
注：\citet{imbens1997estimating} 和 \citet{kitagawa2015test} 讨论了类似结果。可以基于 Kolmogorov--Smirnov statistic 的类似量检验第一条 inequality：
$$
\text{KS}_1 = \max_{y} \Big |
\frac{  \sumn Z_i D_i  1(Y_i \leq y)  }{ \sumn Z_i D_i  } - \frac{  \sumn (1-Z_i) D_i  1(Y_i \leq y)  }{ \sumn (1-Z_i) D_i  }
\Big |.
$$



\paragraph{Example for the IV inequalities}

给出一个所有 IV inequalities 都成立的例子，再给出一个并非所有 IV inequalities 都成立的例子。需要指定 binary variables $(Z,D,Y)$ 的 joint distribution。

 
 
 \paragraph{Violations of the key assumptions}\label{eq::violations-of-IV}

\cref{thm::CACE-identification} 依赖 randomization、monotonicity 和 exclusion restriction。后两个假设即使在 randomized experiments 中也不可检验。当它们被违反时，IV estimator 不再识别 CACE。本题给出下面两个情形，它们是 \citet{angrist1996identification} 中 Propositions 2 and 3 的重述。回忆 $\pi_u = \pr(U=u)$ 且 $\tau_u = E\{  Y(1) - Y(0) \mid U=u \}$，其中 $u=\at, \nt, \cp, \df$。


\begin{theorem}
\label{theorem::wald-violation-assumptions}
(a) 
在 \cref{assume::random} 和 \cref{assume::Mono} 下，但不要求 exclusion restriction，有
$$
\frac{ E(Y\mid Z=1) - E(Y\mid Z=0) }{ E(D\mid Z=1) - E(D \mid Z=0) } - \tau_\cp 
= \frac{  \pi_\at \tau_\at + \pi_\nt \tau_\nt  }{  \pi_\cp } .
$$
 

(b)
在 \cref{assume::random} 和 \cref{assume::ER} 下，但不要求 monotonicity，有
$$
\frac{ E(Y\mid Z=1) - E(Y\mid Z=0) }{ E(D\mid Z=1) - E(D \mid Z=0) } - \tau_\cp 
= \frac{  \pi_\df (\tau_\cp + \tau_\df )  }{   \pi_\cp - \pi_\df   }.
$$
\end{theorem}
 
证明 \cref{theorem::wald-violation-assumptions}。

 
 \paragraph{Problems of other analyses}\label{problem::as-treated-analysis}

在 \cref{sec::disentangle-mixtures} 推导 IV inequalities 的过程中，我们通过识别 latent strata 的比例以及它们 potential outcomes 的 conditional means 来 disentangle mixture distributions。这些结果有助于理解其他看似合理的 analyses 的缺陷。下面回顾三个 estimators，并假设 \cref{assume::random}--\cref{assume::ER} 成立。


\begin{enumerate}
\item 
{\it As-treated analysis} 比较实际接受 treatment 与 control 的 units 的 outcome 均值，得到
$$
\tau_\textsc{at} = E(Y\mid D=1) - E(Y\mid D=0).
$$
证明
$$
\tau_\textsc{at} = \frac{  \pi_\at \mu_\at + \pr(Z=1)\pi_\cp \mu_{1\cp}   }{   \pr(D=1) } 
- \frac{  \pi_\nt \mu_\nt + \pr(Z=0) \pi_\cp \mu_{0\cp}   }{   \pr(D=0) } .
$$


\item 
{\it Per-protocol analysis} 比较 treatment group 与 control group 中遵从所分配 treatment 的 units，得到
$$
\tau_\textsc{pp} = E(Y\mid Z=1, D=1) - E(Y\mid Z=0, D=0).
$$
证明
$$
\tau_\textsc{pp} =\frac{  \pi_\at \mu_\at + \pi_\cp \mu_{1\cp} }{  \pi_\at   + \pi_\cp }
- \frac{  \pi_\nt \mu_\nt + \pi_\cp \mu_{0\cp} }{  \pi_\nt   + \pi_\cp }.
$$


\item 
我们也可能希望在给定 treatment assignment 的条件下，比较接受 treatment 与 control 的 units 的 outcomes，从而得到
\begin{eqnarray*}
\tau_{Z=1} &=& E(Y\mid Z=1, D=1) - E(Y\mid Z=1, D=0),\\ 
\tau_{Z=0} &=& E(Y\mid Z=0, D=1) - E(Y\mid Z=0, D=0).
\end{eqnarray*}
证明它们可化简为
$$
\tau_{Z=1}  = \frac{  \pi_\at \mu_\at + \pi_\cp \mu_{1\cp} }{  \pi_\at   + \pi_\cp } - \mu_\nt,\quad
\tau_{Z=0} = \mu_\at - \frac{  \pi_\nt \mu_\nt + \pi_\cp \mu_{0\cp} }{  \pi_\nt   + \pi_\cp }.
$$
\end{enumerate}
 
 
\paragraph{Bounds on the average causal effect on the whole population}\label{para::bounds-ACE}


基于 \cref{eq::interpret-iv-double} 中的记号，扩展 \cref{sec::disentangle-mixtures} 的讨论。对于 potential outcome $Y(d)$，将 treatment received 对 outcome 的 average causal effect 定义为
$$
\delta = E\{  Y(d=1) - Y(d=0) \},
$$
并由于记号变化，将 $\mu_{zu} $ 的定义改为
$$
m_{du} = E\{  Y(d) \mid U = u\},\quad  (z=0,1;u=\at, \nt, \cp) 
$$  
它们满足
$$
\delta = \sum_{u = \at, \nt, \cp}  \pi_u  (m_{1u} - m_{0u}).
$$

\cref{sec::disentangle-mixtures} 识别了 $\pi_\at$、$\pi_\nt$、$\pi_\cp$、$m_{1\at} = \mu_{1 \at}$、$m_{0\nt} = \mu_{0\nt }$、$m_{1\cp} = \mu_{1\cp}$ 和 $m_{0\cp} = \mu_{0\cp}$。但是，数据不包含关于 $m_{0 \at}$ 和 $m_{1\nt }$ 的任何信息。因此无法识别 $\delta$。当 outcome 有界时，可以给出 $\delta$ 的 bounds。

\begin{theorem}
\label{thm::bound-on-ACE}
在 \cref{assume::Mono}--\cref{assume::ER-double-index} 下，若 outcome 有界于 $[\underline{y}, \overline{y}]$，则有 $\underline{\delta} \leq  \delta  \leq \overline{\delta}$，其中
$$
\underline{\delta} =  \delta'  - \overline{y} \pr(D=1\mid Z=0) + \underline{y} \pr(D=0\mid Z=1)
$$
以及
$$
\overline{\delta} = \delta'  -   \underline{y} \pr(D=1\mid Z=0) +  \overline{y}  \pr(D=0\mid Z=1) 
$$
且 $\delta' = E(DY\mid Z=1) - E(Y-DY\mid Z=0) .$
\end{theorem}

证明 \cref{thm::bound-on-ACE}。
 
 
注：在 binary outcome 的特殊情形下，bounds 可简化为
 $$
\underline{\delta} =   E(DY\mid Z=1) - E( D+Y-DY\mid Z=0)  
$$
以及
$$
\overline{\delta} =  E(DY+1-D\mid Z=1) - E(Y-DY\mid Z=0) .
 $$
 
 
 \paragraph{One-sided noncompliance and statistical inference}\label{para::15-onesidednoncompliance}
 
考虑一个 randomized encouragement design，其中被分配到 control 的 units 无法获得 treatment。对 unit $i$，令 $Z_i$ 为 binary treatment assigned，$D_i$ 为 binary treatment received，$Y_i$ 为感兴趣的 outcome。One-sided noncompliance 指的是
$$
Z_i=0 \Longrightarrow  D_i = 0 \quad (i=1,\ldots, n). 
$$
假设 \cref{assume::random} 成立。
 
 
\begin{enumerate}
[(1)]
\item 
\cref{assume::Mono} 在这个情形下是否成立？
本题中由 $\{ D_i(1), D_i(0)  \}$ 定义的 latent strata 有多少个？如何用 observed data distribution 识别它们的比例？

\item 
陈述 exclusion restriction 的假设。在 exclusion restriction 下，证明 $E\{  Y(z) \mid U = u \}$ 可以由 observed data distributions 识别。
给出所有可能的 $z$ 和 $u$ 取值下的公式。在这个情形中如何识别 CACE？

 

\item
如果对所有 units $i$ 都观测到 pretreatment covariates $X_i$，如何利用 covariate information 提高 CACE 的 estimation efficiency？
 
 
 \item 
 在 \cref{assume::random} 下，exclusion restriction \cref{assume::ER} 有 testable implications，即 one-sided noncompliance 下的 IV inequalities。请陈述这些 IV inequalities。
 
\item
\citet{sommer1991estimating} 给出了如下数据集：
 
 \begin{center}
\begin{tabular}{cccccc} 
\hline 
      &   \multicolumn{2}{c}{$Z=1$}  &  &  \multicolumn{2}{c}{$Z=0$}  \\  \cline{2-3} \cline{5-6} 
      & $D=1$  &  $D=0$           &       &  $D=1$  &  $D=0$\\
 $Y=1$ &   9663    &     2385         &      &     0       &    11514       \\
 $Y=0$ &   12   &     34            &      &     0         &    74           \\  
 \hline 
\end{tabular}
 \end{center}
 
Treatment assigned $Z$ 表示儿童是否被分配到 vitamin A supplement，treatment received indicator $D$ 表示儿童是否实际接受 vitamin A supplement，binary outcome $Y$ 是生存指示变量（survival indicator）。原始 RCT 在 Indonesia 进行。
请重新分析这组数据。
 \end{enumerate}
 
 
注：\citet{bloom1984accounting} 最早讨论 one-sided noncompliance，并提出 estimator $\hat\tau_\cp = \hat\tau_Y / \hat\tau_D$；它有时称为 Bloom estimator。\citet{bloom1984accounting} 的记号与本章不同。

 
 
 
 
  \paragraph{One-sided noncompliance with partial adherence}
\label{para::one-sided-partial-adherence}
 
\citet[][Table 3]{sanders2021incorporating} 报告了如下 RCT 数据；该 RCT 旨在估计戒烟干预（smoking cessation interventions）在 psychotic disorders 个体中的有效性（efficacy）。

\begin{center}
\begin{tabular}{cccc}
\hline  
分配组 & treatment received & 样本量 & \# positive outcomes \\
\hline  
Control & None& 151 &25 \\
Treatment& None& 35 & 7 \\
Treatment& Partial& 42& 17 \\
Treatment& Full& 70& 40\\
\hline 
\end{tabular}
\end{center}

Treatment received 有三个层级：``full'' treatment 对应参加全部 $8$ 次 treatment sessions，``partial'' 对应参加 $5$ 到 $7$ 次 sessions，``none'' 对应参加少于 $5$ 次 sessions。Outcome 定义为三个月时相对于 baseline 的 smoking reduction 是否达到 $50\%$ 或以上的 binary indicator。


在本题中，treatment assignment $Z$ 是 binary，但 treatment received $D$ 对 ``none''、``partial'' 和 ``full'' 分别取三个值 $0, 0.5, 1$。三水平的 $D$ 会带来复杂性，但在 control assignment 下它只能为 $0$。本题中有多少个 latent strata $U = \{ D(1), D(0) \} $？能否识别它们的比例？

如何将 exclusion restriction 推广到本题？哪些 causal effects 可以作为感兴趣的对象？它们能否被识别？

请基于上述问题分析这组数据。



  

 \paragraph{Recommended reading}

\citet{balke1997bounds} 推导了更一般的 IV inequalities。
